\subsection{Getting Corvus}
To get corvus, go to 
\url{https://github.com/times-software/Corvus/releases/latest}
and download the archived .zip or .tgz file.

\subsubsection{Requirements}
Corvus is portable and runs on macOS, Linux, and Windos. 
To install corvus, you will need python. Versions > 3.12 and < 3.14 are supported. Earlier versions may work as well. We also strongly suggest that you use
an environment manager such as conda to install corvus. Our suggestion is to
use the free miniforge package, which can be found at \url{https://conda-forge.org/download/}

Python's venv environment management package can also be used.

To actually run corvus, you will need at least FEFF10. To get FEFF10, go to \url{https://times-software.github.io/feff10/} to register. Registration is free,
and is only used for communication about bugs and or updates, and for usage statistics. After registration, you will receive an email with directions for 
download and installation.

\subsection{Installation}

\subsubsection{Simple Installation using install script}
Once you have a version of Corvus from the github website, you can follow the following simple directions to install. Note that in the below "version" should be
replaced with the version that you downloaded, e.g., Corvus-3.2.0-01.zip.

\paragraph{MacOS} Unzip the downloaded file. Open a terminal and navigate to the file. Run the install.sh script.
\begin{verbatim}
cd Downloads
unzip Corvus-version.zip
cd Corvus-version/
source install.sh
\end{verbatim}

\paragraph{Windows} Unzip the download file. Open a terminal and navigate to the directory. Run the installer.bat script.
\begin{verbatim}
cd Downloads
unzip Corvus-version.zip
cd Corvus-version\Corvus-version
install.bat
\end{verbatim}

\subsubsection{Step-by-step instructions with conda}
In MacOS, just open a terminal, and it will already have conda enabled if you have it installed. In windows, you will need to open a conda terminal, (Miniforge command window if you installed miniforge conda).
\begin{enumerate}
\item Download corvus from the github page.
\item unzip the folder.
\item In the terminal, navigate to the folder
\item Create a conda environment for corvus
\item Install corvus with pip.
\end{enumerate}
\begin{verbatim}
cd Downloads/Corvus-version/
conda create -name your_corvus_env python=">=3.12,<3.14"
conda activate your_corvus_env
pip install .
\end{verbatim}
